\documentclass[11pt]{article}
\usepackage[margin=1in]{geometry}
\usepackage{amsmath,amssymb,amsthm}
\usepackage{enumitem}
\usepackage{booktabs}
\usepackage{xcolor}
\usepackage{tcolorbox}
\usepackage{hyperref}
\tcbuselibrary{skins,breakable}

\newtheorem{definition}{Definition}
\newtheorem{lemma}{Lemma}
\newtheorem{remark}{Remark}

\newtcolorbox{correction}[1][]{%
  colback=red!4,colframe=red!55!black,title=Correction,fonttitle=\bfseries,
  breakable,#1}
\newtcolorbox{clarify}[1][]{%
  colback=blue!4,colframe=blue!45!black,title=Clarification,fonttitle=\bfseries,
  breakable,#1}
\newtcolorbox{nuance}[1][]{%
  colback=yellow!8,colframe=orange!70!black,title=Implementation nuance (not a paper error),fonttitle=\bfseries,
  breakable,#1}

\title{Ambient, Contact, and Boundary Graphs for Rauzy Fractals:\\
a Working Reference with Corrections}
\author{Reverification branch notes\\
\small Restated independently while reimplementing and testing the constructions of:\\
\small A.~Loridant, J.~Thuswaldner, Y.~Zhang, \emph{Neighbors of self-affine tiles and Rauzy fractals},\\
\small arXiv:2511.16442 (November 2025)}
\date{}

\begin{document}
\maketitle

\begin{abstract}
This note restates, in our own words and notation, the parts of the cited paper
needed to construct the ambient, contact, and self-replicating boundary graphs
of a Pisot substitution, and to grow the boundary graph from the contact graph
via the ``C-corona'' procedure (their Algorithm 2). It exists because an
independent, from-scratch reimplementation of these constructions surfaced a
handful of concrete transcription errors in the published worked examples, one
genuinely easy-to-miss notational ambiguity that caused a real implementation
bug, and one point where our own simplification (not the paper) was the source
of an apparent discrepancy. All three are catalogued explicitly in
Section~\ref{sec:errata}, each tied to how it was actually found: by
independent computation from first principles, not by re-reading the source
more carefully. Every claim below was checked computationally against the
paper's own two worked examples ($\sigma_1$, $\sigma_2$) before being trusted;
Section~\ref{sec:examples} gives the corrected data.
\end{abstract}

\section{Setup}

Let $\sigma$ be a substitution on a finite alphabet $\mathcal{A}=\{1,\dots,d\}$,
primitive, with an irreducible characteristic polynomial for its incidence
matrix $\mathbf{M}$ (so $\mathbf{M}$ has a simple dominant eigenvalue $\beta>1$,
the Pisot number associated to $\sigma$; the remaining eigenvalues are strictly
smaller in modulus). Write $\mathbf{u}$ for the right Perron eigenvector,
normalized to sum $1$ (letter frequencies), and $\mathbf{v}$ for the left
Perron eigenvector, normalized however is convenient (any positive scalar
multiple gives an equivalent construction below).

For a word $w$ over $\mathcal{A}$, $l(w)\in\mathbb{Z}^d$ denotes its
abelianization (the vector of letter counts).

\begin{definition}[Prefix--suffix graph]
The prefix--suffix data $\mathcal{P}$ is the set of triples $(p,i,s)$, with
$p,s$ words over $\mathcal{A}$ and $i\in\mathcal{A}$, such that
$\sigma(j)=p\,i\,s$ for some letter $j$. We record this as a directed edge
$i\to j$, labelled $(p,i,s)$: informally, ``$j$ is a parent of $i$, and $i$
occurs in $\sigma(j)$ preceded by $p$ and followed by $s$.''
\end{definition}

\begin{definition}[Stepped hyperplane $H_\sigma$]
$$H_\sigma = \{\,[x,i]\in\mathbb{Z}^d\times\mathcal{A} \;:\; 0\le \langle x,\mathbf{v}\rangle < \langle e_i,\mathbf{v}\rangle\,\}.$$
\end{definition}

This is the coordinatization used throughout: a face of the stepped surface is
identified by an integer vector $x$ (its position) together with a colour
$i\in\mathcal{A}$ (which coordinate direction it omits). Testing membership in
$H_\sigma$ therefore reduces to a single linear inequality in $\langle
x,\mathbf{v}\rangle$; because $\mathbf{v}$ is generally irrational, we found it
important to make this test genuinely exact rather than floating-point-with-a-tolerance
(Remark~\ref{rem:exact} below).

\begin{definition}[Dual substitution $\sigma^*$]
For $[x,i]\in H_\sigma$,
$$\sigma^*\{[x,i]\} = \{\,[\mathbf{M}^{-1}(x+l(p)),\,j]\;:\;i\to j \text{ an edge labelled } (p,i,s)\,\}.$$
\end{definition}

\begin{lemma}[Invariance]
$\sigma^*$ maps $H_\sigma$ into $H_\sigma$, and the images arising from
distinct source faces are pairwise disjoint; iterating $\sigma^*$ partitions
$H_\sigma$ into ever finer pieces.
\end{lemma}

This invariance is the reason $\sigma^*$ never needs an ad hoc ``fold back into
range'' step: boundedness is a theorem about the arithmetic, not something
enforced after the fact. We flag this because a prior attempt at this
implementation used a hand-rolled folding operation instead of routing
everything through $\sigma^*$, and that fold turned out not to commute with
the matrix recursion under iteration -- the single largest source of wasted
effort before this reimplementation.

\section{Graphs describing the tile boundary}

Fix the set
$$\mathfrak{D} = \{\,[i,x,j]\in\mathcal{A}\times H_\sigma \;:\; x=0 \Rightarrow i<j\,\}$$
(here ``$[i,x,j]\in\mathcal{A}\times H_\sigma$'' means $i$ ranges freely over
$\mathcal{A}$ while $[x,j]\in H_\sigma$; the ordering condition only bites when
$x=0$, and exists purely to avoid double-counting $[i,0,j]$ and $[j,0,i]$,
which denote the same object).

\begin{definition}[Ambient graph]\label{def:ambient}
The nodes of the ambient graph $G_\mathfrak{D}$ are the elements of
$\mathfrak{D}$. There is an edge $[i,x,j]\to[i',x',j']$ whenever there exist
$(p_1,i,s_1),(q_1,j,t_1)\in\mathcal{P}$ with either
\begin{align*}
\text{(type 1)}\quad &\sigma(i')=p_1\,i\,s_1,\ \ \sigma(j')=q_1\,j\,t_1,\ \ \mathbf{M}x' = x+l(q_1)-l(p_1),\\
\text{(type 2)}\quad &\sigma(j')=p_1\,i\,s_1,\ \ \sigma(i')=q_1\,j\,t_1,\ -\mathbf{M}x' = x+l(q_1)-l(p_1).
\end{align*}
\end{definition}

\begin{definition}[Self-replicating boundary graph, contact graph]
Write $\mathrm{Red}(G)$ for the largest subgraph of $G$ with no node lacking
an outgoing edge (remove zero-outdegree nodes repeatedly until none remain).
The self-replicating boundary graph $G_B$ is the largest subgraph of
$G_\mathfrak{D}$ in which every node lies on some walk that eventually enters
a cycle. Separately, the (pre-)contact graph $G_P$ is the largest subgraph of
$G_\mathfrak{D}$ in which \emph{every} walk eventually reaches a
distinguished finite seed set $\mathfrak{D}_{\mathrm{cont}}$ (the ``contact
set,'' geometrically: pairs of faces $[0,i]$, $[x,j]$ whose intersection is
exactly $(d-2)$-dimensional); the contact graph is $G_C=\mathrm{Red}(G_P)$.
$G_P$ -- and hence $G_C$ -- can be computed by starting from
$\mathfrak{D}_{\mathrm{cont}}$ and repeatedly adjoining any node with an edge
into the current set, which provably terminates.
\end{definition}

\subsection{The signed (``simple'') variants}

Because $\mathfrak{D}$ only keeps one of $[i,x,j]$ and its mirror image for
each geometric object, it is convenient to also work in the fully signed
space, using
$$-[i,x,j] := [j,-x,i]$$
(swap the two colours and negate the offset). We will call this
\textbf{triple-level negation}, written $\ominus$, to keep it visually distinct
from a second, unrelated negation introduced just below.

\begin{definition}[Simple ambient graph]
Nodes are $\pm(\mathcal{A}\times H_\sigma)\setminus\{[i,0,i]:i\in\mathcal{A}\}$.
Edges use \emph{only} the type-1 rule of Definition~\ref{def:ambient} (type 2
is not needed separately: once both a node and its $\ominus$-image are
present, a type-1 edge out of one plays the role type 2 played in the
restricted graph). $\hat G_B$, $\hat G_P$, $\hat G_C=\mathrm{Red}(\hat G_P)$
are defined exactly as before, now inside this larger signed space. Every node
of $\mathfrak{D}$ gives rise to itself and its $\ominus$-image in the signed
graphs, and conversely every signed node arises this way from some node of
$\mathfrak{D}$ (possibly after $\ominus$).
\end{definition}

\begin{correction}
\textbf{Verifying the type-1/type-2 correspondence.} It is tempting to guess
that a type-2 edge $[i,x,j]\to[i',x',j']$ in the restricted graph simply
corresponds to a type-1 edge from $\ominus[i,x,j]$ to $\ominus[i',x',j']$.
Working the equations through carefully (matching which prefix attaches to
which parent letter) shows this guess is \emph{false}: the correct
correspondence is a type-1 edge from $[i,x,j]$ itself (unchanged) to
$\ominus[i',x',j'] = [j',-x',i']$. We record this because the wrong guess is
the more natural-looking one, and it would silently produce an internally
self-consistent but semantically incorrect implementation -- self-consistency
under an independent forward/backward differential test does \emph{not} rule
this class of error out, since both directions can share the same mistaken
correspondence.
\end{correction}

\subsection{Connections and the C-corona}

\begin{definition}[Connections]
An element $[i_1,x,i_2]\in\pm(\mathcal{A}\times H_\sigma)$ is regarded as a
(symmetric) connection between $[y,i_1]$ and $[y+x,i_2]$, for any
$y\in\mathbb{Z}^d$. Let $C=\mathrm{nodes}(G_C)$ be the contact set. Given
$[x_1,i_1],[x_2,i_2]\in H_\sigma$, we say they are joined by $r$ subsequent
$\pm C$-connections, written $[x_1,i_1]\sim_r[x_2,i_2]$, if there is a chain
$[y_0,j_0]=[x_1,i_1],\dots,[y_r,j_r]=[x_2,i_2]$ with
$[j_k,\,y_{k+1}-y_k,\,j_{k+1}]\in\pm C\cup\{[i,0,i]:i\in\mathcal{A}\}$ for each
step.
\end{definition}

\begin{clarify}\label{clar:negation}
\textbf{Two different negations, one overloaded symbol.} The relation
$\sim_r$ is stated above for $[x_1,i_1],[x_2,i_2]\in H_\sigma$. The source
material extends it to the case where both endpoints lie in $-H_\sigma$ by
declaring $[x_1,i_1]\sim_r[x_2,i_2]$ to hold exactly when
$[-x_1,i_1]\sim_r[-x_2,i_2]$ does. Notice what is negated: only the vector
$x_1$, \emph{not} the colour $i_1$. This is a second, genuinely different
operation from the $\ominus$ used to build the signed node space one
subsection above -- it acts on a \emph{pair} $[x,i]$, keeps the letter fixed,
and flips only the vector:
$$\boxminus[x,i] := [-x,i], \qquad -H_\sigma := \{[x,i] : \boxminus[x,i]\in H_\sigma\}.$$
We introduce the boxed-minus symbol here specifically because the source
material uses the same ``$-$'' for both $\ominus$ (triple, swap-and-negate,
used for node validity in $\pm(\mathcal{A}\times H_\sigma)$) and $\boxminus$
(pair, negate-only, used inside the definition of $\sim_r$ and hence of the
C-corona below). Conflating the two is exactly the bug we hit: filtering
corona candidates by whether the whole triple was $\ominus$-valid (which is
the right test for node membership) instead of by whether the relevant
\emph{pair} was in $H_\sigma\cup\boxminus H_\sigma$ (the right test for
whether a connection endpoint is well-formed) let in a large number of
spurious candidates and made the corona iteration diverge instead of converge.
Both endpoints of every $\pm C$-hop -- the one you are extending from and the
one you land on -- need the $\boxminus$-style check, not the $\ominus$-style
one.
\end{clarify}

\begin{definition}[C-corona]
For a subgraph $\hat G$ of the simple ambient graph,
$$\mathrm{C\text{-}Corona}(\hat G) = \Big\{[i_1,x,i_2]\in\pm(\mathcal{A}\times H_\sigma)\setminus\{[i,0,i]\} \;:\; \exists\,[i_1,y,j]\in\hat G \text{ with } [y,j]\sim_1[x,i_2]\Big\}.$$
Concretely: for each node $[i_1,y,j]$ already present, and each connector
$[j,\delta,i_2]\in\pm C\cup\{[k,0,k]\}$ whose starting colour matches $j$
(after the $\boxminus$-validity check of Clarification~\ref{clar:negation} on
both $[y,j]$ and $[y+\delta,i_2]$), adjoin $[i_1,\,y+\delta,\,i_2]$.
\end{definition}

\subsection{Algorithm 2}

\begin{quote}
$A_1 \leftarrow \hat G_C$. Repeat $A_{p} \leftarrow \mathrm{Red}(\mathrm{C\text{-}Corona}(A_{p-1}))$ until $A_p = A_{p-1}$; the fixed point is $\hat G_B$.
\end{quote}

\begin{nuance}\label{nuance:2826}
\textbf{$\hat G_C$ is not simply $C\cup(-C)$ handled naively.} Building the
starting set as the raw union of $C$ (the restricted contact set) with its
$\ominus$-images is the natural first attempt, and it does produce a graph
with the right nodes -- but applying $\mathrm{Red}$ to that raw union \emph{in
isolation}, using only edges among those nodes, does not shrink it, even when
(as for the example below) the truly-correct fixed-point size is smaller.
This is not a defect in the raw union; it simply means node membership in
$\hat G_C$, done properly, is a global property of the whole reachability
structure, not something you can shortcut by looking only at the small
seed set before the corona iteration has had a chance to explore outward and
prune back. Running the full iteration (rather than trying to pre-trim the
seed) resolves this correctly on its own -- we mention it only so a future
implementer does not mistake a raw, unreduced union for a discrepancy with
the literature.
\end{nuance}

\begin{remark}[On exact arithmetic]\label{rem:exact}
Every quantity above that looks like it needs a numerical tolerance --
membership in $H_\sigma$, chiefly -- can instead be decided exactly. Represent
$\beta$ as a specific real root of the (irreducible) characteristic
polynomial, and $\mathbf{u},\mathbf{v}$ as elements of the number field
$\mathbb{Q}(\beta)$, i.e.\ as rational-coefficient polynomials in $\beta$ of
degree less than $d$, reduced modulo the minimal polynomial. Any such element
is exactly zero iff its reduced representative has all-zero coefficients --
decidable purely over $\mathbb{Q}$, with no floating point at all. Only once a
quantity is known to be nonzero does one need a numeric value, purely to read
off its sign; at that point there is no remaining ambiguity to get wrong. We
found this worth the extra machinery: it turns a family of ``is this exactly
on the boundary'' bugs into a non-issue rather than a source of
epsilon-tuning.
\end{remark}

\section{Worked examples}\label{sec:examples}

Let $\sigma_1: 1\mapsto 1112,\ 2\mapsto 113,\ 3\mapsto 1$ and
$\sigma_2: 1\mapsto 112,\ 2\mapsto 1113,\ 3\mapsto 1$, both cubic
(characteristic polynomials $x^3-3x^2-2x-1$ and $x^3-2x^2-x-1$ respectively).

\begin{correction}
\textbf{Two transcription errors in the published contact sets.} Computing
$\mathfrak{D}_{\mathrm{cont}}$ directly from the geometric definition (which
face pairs intersect in an exactly $(d-2)$-dimensional cube), independently of
the published lists, reproduces the published nine-element sets for both
$\sigma_1$ and $\sigma_2$ to 8/9 and 7/9 respectively -- and in every
mismatched slot, the independently-computed replacement differs from the
published entry only in a single small, easy-to-typo way:
\begin{itemize}[leftmargin=1.5em]
\item $\sigma_1$: published $[2,e_3,1]$ should read $[2,e_3,2]$ (a same-colour
self-contact, structurally parallel to the $[1,e_2,1]$ and $[1,e_3,1]$ entries
already in the same list -- not a $2$-$1$ contact at all).
\item $\sigma_2$: published $[2,e_3,1]$ should likewise read $[2,e_3,2]$; and
published $[1,e_2,2]$ should read $[1,e_2-e_1,2]$ (an apparently dropped
$-e_1$ term).
\end{itemize}
Corrected sets:
$$\mathfrak{D}_{\mathrm{cont}}(\sigma_1)=\{[1,e_2,1],[1,e_3,1],[1,0,2],[1,0,3],[2,e_1{-}e_2,1],[2,e_3,2],[2,0,3],[3,e_1{-}e_3,1],[3,e_2{-}e_3,2]\}$$
$$\mathfrak{D}_{\mathrm{cont}}(\sigma_2)=\{[1,e_2{-}e_1,2],[1,e_3,1],[1,0,2],[1,0,3],[2,e_3,2],[2,e_1,2],[2,0,3],[3,e_1{-}e_3,1],[3,e_2{-}e_3,2]\}$$
\end{correction}

\begin{table}[h]
\centering
\begin{tabular}{@{}lccccc@{}}
\toprule
& $|\mathfrak{D}_{\mathrm{cont}}|$ & $|C|=|G_C|$ & $|\hat G_C|$ & $|\hat G_B|$ & growth? \\
\midrule
$\sigma_1$ & 9 & 14 & 26 & 26 & none (matches source claim) \\
$\sigma_2$ & 9 & 15 & 30 & 50 & yes \\
\bottomrule
\end{tabular}
\caption{Corrected/reproduced example data. All four columns for $\sigma_1$
were independently reproduced exactly from a from-scratch implementation,
including the source's claim that the boundary graph and contact graph
coincide for $\sigma_1$. $\sigma_2$'s $|\hat G_B|=50$ has no published number
to check against (the source shows it only as a figure); it is reported here,
not verified against the literature.}
\end{table}

\section{Errata and clarifications, at a glance}\label{sec:errata}

\begin{enumerate}[leftmargin=1.5em]
\item \textbf{Genuine transcription errors} (Section~\ref{sec:examples}):
two single-character-scale slips in the published $\mathfrak{D}_{\mathrm{cont}}$
lists, found by computing the sets from the geometric definition rather than
trusting the tables, and confirmed by the fact that each replacement
restores an otherwise-missing structural symmetry (a same-colour self-contact
entry, or a plausible dropped term).
\item \textbf{Genuine notational ambiguity} (Clarification~\ref{clar:negation}):
the source overloads ``$-$'' for two different operations -- swap-and-negate
on triples (for node validity) versus negate-only on pairs (for the $\sim_r$
connection relation and hence the C-corona). The text is not, in the end,
self-contradictory -- the pair-level convention is recoverable from a close
reading of how $\sim_r$ extends to $-H_\sigma$ -- but the two are easy to
conflate, and doing so produces code that passes an internal
forward/backward consistency check while still being wrong, because both
directions inherit the same mistake.
\item \textbf{Our own implementation shortcut, not a paper error}
(Nuance~\ref{nuance:2826}): building $\hat G_C$ as a raw union $C\cup\ominus C$
and expecting it to already equal the correct fixed-point size is not
justified; the full Algorithm~2 iteration is what actually determines
membership, and running it resolves the apparent discrepancy on its own.
\end{enumerate}

\end{document}
